\documentclass[12pt,a4paper]{article}
\usepackage[utf8]{inputenc}
\usepackage[T1]{fontenc}
\usepackage[french]{babel}
\usepackage{geometry}
\usepackage{listings}
\usepackage{xcolor}
\usepackage{titlesec}
\usepackage{fancyhdr}
\usepackage{booktabs}
\usepackage{graphicx}
\usepackage{array}
\usepackage{mdframed}
\usepackage{enumitem}
\usepackage{float}
\usepackage{tabularx}
\usepackage[hidelinks]{hyperref}
\usepackage{hyperref}

\geometry{top=2.5cm, bottom=2.5cm, left=2.5cm, right=2.5cm}

% ── Couleurs
\definecolor{terraformPurple}{RGB}{94, 56, 163}
\definecolor{pulumiBlue}{RGB}{18, 97, 220}
\definecolor{cicdGreen}{RGB}{0, 128, 80}
\definecolor{codeGray}{RGB}{245, 245, 245}
\definecolor{codeBorder}{RGB}{200, 200, 200}
\definecolor{successGreen}{RGB}{34, 139, 34}
\definecolor{terminalBg}{RGB}{22, 22, 30}
\definecolor{terminalFg}{RGB}{180, 255, 180}
\definecolor{keywordBlue}{RGB}{0, 100, 200}
\definecolor{warningOrange}{RGB}{200, 100, 0}
\definecolor{sectionTf}{RGB}{94, 56, 163}
\definecolor{sectionPu}{RGB}{18, 97, 220}

% ── Styles listings
\lstdefinelanguage{HCL}{
  keywords={terraform,provider,resource,variable,output,required_providers,source,version,default,description,type,string,number,bool,true,false,depends_on,keep_locally,build,context,dockerfile,internal,external,name,image,ports,env},
  keywordstyle=\color{terraformPurple}\bfseries,
  sensitive=true,
  comment=[l]{\#},
  commentstyle=\color{gray}\itshape,
  stringstyle=\color{orange!80!black},
  morestring=[b]",
}

\lstdefinelanguage{YAML}{
  keywords={name,on,push,branches,jobs,runs-on,steps,uses,with,run,if},
  keywordstyle=\color{cicdGreen}\bfseries,
  sensitive=true,
  comment=[l]{\#},
  commentstyle=\color{gray}\itshape,
  stringstyle=\color{orange!80!black},
  morestring=[b]",
  morestring=[b]',
}

\lstdefinelanguage{TypeScript}{
  keywords={import,export,const,let,var,new,from,async,await,true,false,return},
  keywordstyle=\color{pulumiBlue}\bfseries,
  sensitive=true,
  comment=[l]{//},
  morecomment=[s]{/*}{*/},
  commentstyle=\color{gray}\itshape,
  stringstyle=\color{orange!80!black},
  morestring=[b]",
  morestring=[b]',
  morestring=[b]`
}

\lstdefinestyle{hclStyle}{
  language=HCL,
  backgroundcolor=\color{terraformPurple!5},
  basicstyle=\ttfamily\footnotesize,
  frame=single,
  rulecolor=\color{terraformPurple!40},
  breaklines=true,
  showstringspaces=false,
  numbers=left,
  numberstyle=\tiny\color{gray},
  numbersep=8pt,
  tabsize=2,
  captionpos=b
}

\lstdefinestyle{yamlStyle}{
  language=YAML,
  backgroundcolor=\color{cicdGreen!5},
  basicstyle=\ttfamily\footnotesize,
  frame=single,
  rulecolor=\color{cicdGreen!40},
  breaklines=true,
  showstringspaces=false,
  numbers=left,
  numberstyle=\tiny\color{gray},
  numbersep=8pt,
  tabsize=2,
  captionpos=b
}

\lstdefinestyle{tsStyle}{
  language=TypeScript,
  backgroundcolor=\color{pulumiBlue!5},
  basicstyle=\ttfamily\footnotesize,
  frame=single,
  rulecolor=\color{pulumiBlue!40},
  breaklines=true,
  showstringspaces=false,
  numbers=left,
  numberstyle=\tiny\color{gray},
  numbersep=8pt,
  tabsize=2,
  captionpos=b
}

\lstdefinestyle{termStyle}{
  backgroundcolor=\color{terminalBg},
  basicstyle=\ttfamily\footnotesize\color{terminalFg},
  frame=single,
  rulecolor=\color{black},
  breaklines=true,
  showstringspaces=false,
  numbers=none,
  tabsize=2,
  captionpos=b
}

\lstdefinestyle{dockerStyle}{
  language=bash,
  backgroundcolor=\color{gray!10},
  basicstyle=\ttfamily\footnotesize,
  frame=single,
  rulecolor=\color{gray!50},
  breaklines=true,
  showstringspaces=false,
  numbers=left,
  numberstyle=\tiny\color{gray},
  tabsize=2,
  captionpos=b
}

% ── En-tête / pied de page
\pagestyle{fancy}
\fancyhf{}
\rhead{\textcolor{gray}{\small TP IaC --- Terraform \& Pulumi}}
\lhead{\textcolor{gray}{\small DevOps 2026}}
\rfoot{\thepage}
\renewcommand{\headrulewidth}{0.4pt}

% ── Boîtes
\newmdenv[linecolor=terraformPurple,backgroundcolor=terraformPurple!5,
  linewidth=2pt,roundcorner=4pt,
  innerleftmargin=10pt,innerrightmargin=10pt,
  innertopmargin=8pt,innerbottommargin=8pt]{tfbox}

\newmdenv[linecolor=pulumiBlue,backgroundcolor=pulumiBlue!5,
  linewidth=2pt,roundcorner=4pt,
  innerleftmargin=10pt,innerrightmargin=10pt,
  innertopmargin=8pt,innerbottommargin=8pt]{pubox}

\newmdenv[linecolor=cicdGreen,backgroundcolor=cicdGreen!5,
  linewidth=2pt,roundcorner=4pt,
  innerleftmargin=10pt,innerrightmargin=10pt,
  innertopmargin=8pt,innerbottommargin=8pt]{cicdbox}

\newmdenv[linecolor=successGreen,backgroundcolor=successGreen!5,
  linewidth=2pt,roundcorner=4pt,
  innerleftmargin=10pt,innerrightmargin=10pt,
  innertopmargin=8pt,innerbottommargin=8pt]{successbox}

\newmdenv[linecolor=warningOrange,backgroundcolor=warningOrange!5,
  linewidth=2pt,roundcorner=4pt,
  innerleftmargin=10pt,innerrightmargin=10pt,
  innertopmargin=8pt,innerbottommargin=8pt]{warnbox}

\titleformat{\section}{\large\bfseries}{}{0em}{\thesection\quad}[\titlerule]
\titleformat{\subsection}{\normalsize\bfseries}{}{0em}{\thesubsection\quad}

% ════════════════════════════════════════════
\begin{document}

% ── Page de titre
\begin{titlepage}
  \centering
  \vspace*{1.5cm}

  {\Huge\bfseries Travaux Pratiques\\[0.5em]
  Infrastructure as Code (IaC)}\\[0.5em]
  {\Large\color{gray} Pipelines CI/CD \& Automatisation DevOps}

  \vspace{0.8cm}
  \rule{\linewidth}{1.5pt}
  \vspace{0.5cm}

  \begin{tabular}{ccc}
    \Large\textcolor{terraformPurple}{\textbf{Partie I}} &
    \Large\textcolor{cicdGreen}{\textbf{CI/CD}} &
    \Large\textcolor{pulumiBlue}{\textbf{Partie II}} \\[0.3em]
    \textcolor{terraformPurple}{IaC avec Terraform} &
    \textcolor{cicdGreen}{Pipeline GitHub Actions} &
    \textcolor{pulumiBlue}{IaC avec Pulumi (TypeScript)} \\[0.2em]
    Docker + PostgreSQL + Nginx &
    Automatisation du DLC &
    Docker + LocalStack S3 \\
  \end{tabular}

  \vspace{0.5cm}
  \rule{\linewidth}{0.5pt}
  \vspace{1.5cm}

  \begin{tabular}{ll}
    \textbf{Module :}   & DevOps --- Infrastructure as Code \\[0.4em]
    \textbf{Outils :}   & Terraform CLI, Pulumi CLI, Docker, GitHub Actions \\[0.4em]
    \textbf{Langages :} & HCL (Terraform), TypeScript (Pulumi), YAML (CI/CD) \\[0.4em]
    \textbf{Durée :}    & 6 heures \\[0.4em]
    \textbf{Date :}     & Avril 2026 \\
  \end{tabular}

  \vfill
  {\small\color{gray} Infrastructure 100\% Locale --- Aucune dépendance cloud réelle}
\end{titlepage}

\tableofcontents
\newpage

% ════════════════════════════════════════════
\section{Introduction Générale}

Ce TP couvre deux approches complémentaires de l'\textbf{Infrastructure as Code (IaC)},
ainsi que leur intégration dans un \textbf{pipeline CI/CD} automatisé :

\begin{center}
\begin{tabular}{>{\bfseries}p{3cm} p{4cm} p{4cm} p{3cm}}
\toprule
 & \textbf{\textcolor{terraformPurple}{Partie I}} &
   \textbf{\textcolor{cicdGreen}{CI/CD}} &
   \textbf{\textcolor{pulumiBlue}{Partie II}} \\
\midrule
Outil     & Terraform CLI        & GitHub Actions    & Pulumi CLI \\
Langage   & HCL (déclaratif)     & YAML              & TypeScript \\
Infra     & Docker local         & Runner Ubuntu     & Docker + LocalStack \\
Objectif  & Maîtriser le DLC     & Automatiser l'IaC & Programmer l'infra \\
\bottomrule
\end{tabular}
\end{center}

\vspace{0.5em}
Les trois parties partagent le même objectif : déployer une \textbf{base de données PostgreSQL}
et une \textbf{application web Nginx} de façon reproductible, automatisée et versionnée.

% ════════════════════════════════════════════
% PARTIE I — TERRAFORM
% ════════════════════════════════════════════
\newpage
\begin{center}
  \rule{\linewidth}{2pt}\\[0.5em]
  {\LARGE\bfseries\textcolor{terraformPurple}{PARTIE I}}\\[0.3em]
  {\Large Infrastructure as Code Locale avec Terraform}\\[0.3em]
  \rule{\linewidth}{2pt}
\end{center}

\section{\textcolor{terraformPurple}{Présentation de Terraform}}

Terraform est un outil IaC open-source développé par HashiCorp. Il utilise
un langage déclaratif appelé \textbf{HCL (HashiCorp Configuration Language)}
pour décrire l'infrastructure souhaitée. Terraform calcule ensuite les actions
nécessaires pour atteindre cet état cible.

\begin{tfbox}
\textbf{Objectifs de la Partie I :}
\begin{itemize}[leftmargin=1.5em]
  \item Créer la structure de fichiers Terraform (\texttt{main.tf}, \texttt{variables.tf}, \texttt{outputs.tf})
  \item Déployer PostgreSQL et Nginx via le provider Docker
  \item Exécuter le cycle de vie complet : \texttt{init} \textrightarrow{} \texttt{plan} \textrightarrow{} \texttt{apply} \textrightarrow{} \texttt{destroy}
  \item Valider l'accès à l'application sur \texttt{http://localhost:8080}
\end{itemize}
\end{tfbox}

\section{\textcolor{terraformPurple}{Structure du Projet Terraform}}

\begin{lstlisting}[style=termStyle, caption=Structure du répertoire tp-iac-local/]
tp-iac-local/
├── main.tf            # Ressources et provider Docker
├── variables.tf       # Parametres configurables
├── outputs.tf         # Informations de sortie
├── Dockerfile_app     # Blueprint image Nginx
├── index.html         # Page HTML servie par Nginx
├── nginx.conf         # Configuration du serveur Nginx
└── .gitignore         # Exclusions Git
\end{lstlisting}

\begin{warnbox}
\textbf{Pourquoi 3 fichiers pour le conteneur Nginx ?}
Sur Windows, la commande \texttt{RUN echo} dans un Dockerfile génère des problèmes
d'encodage (BOM/CRLF). La solution adoptée est d'utiliser \texttt{COPY} pour
injecter des fichiers créés localement --- approche plus robuste et portable.
\end{warnbox}

\section{\textcolor{terraformPurple}{Fichiers de Configuration}}

\subsection{Dockerfile\_app}

Image légère basée sur Alpine avec Nginx. La configuration est injectée via
\texttt{COPY} pour éviter les problèmes d'encodage Windows :

\begin{lstlisting}[style=dockerStyle, caption=Dockerfile\_app]
FROM alpine:latest

# Installation de Nginx
RUN apk update && \
    apk add nginx && \
    rm -rf /var/cache/apk/*

# Creation des dossiers requis
RUN mkdir -p /var/www/localhost/htdocs && \
    mkdir -p /run/nginx

# Injection de la page HTML et de la config Nginx
COPY index.html /var/www/localhost/htdocs/index.html
COPY nginx.conf /etc/nginx/http.d/default.conf

EXPOSE 80

CMD ["nginx", "-g", "daemon off;"]
\end{lstlisting}

\subsection{index.html}

\begin{lstlisting}[style=dockerStyle, caption=index.html]
<h1>Application Deployed via Terraform IaC!</h1>
\end{lstlisting}

\subsection{nginx.conf}

La configuration par défaut d'Alpine retourne volontairement 404. Ce fichier
la remplace pour servir correctement les fichiers statiques :

\begin{lstlisting}[style=dockerStyle, caption=nginx.conf]
server {
    listen 80;
    root /var/www/localhost/htdocs;
    index index.html;
    location / {
        try_files $uri $uri/ =404;
    }
}
\end{lstlisting}

\subsection{variables.tf}

\begin{lstlisting}[style=hclStyle, caption=variables.tf]
# --- Variables de Base de Donnees (PostgreSQL) ---
variable "db_name" {
  description = "Nom de la base de donnees PostgreSQL."
  type        = string
  default     = "devops_db"
}

variable "db_user" {
  description = "Nom d'utilisateur PostgreSQL."
  type        = string
  default     = "devops_user"
}

variable "db_password" {
  description = "Mot de passe PostgreSQL."
  type        = string
  default     = "strongpassword123"
}

# --- Variables d'Application ---
variable "app_port_external" {
  description = "Port externe pour l'application web."
  type        = number
  default     = 8080
}
\end{lstlisting}

\subsection{main.tf}

\begin{lstlisting}[style=hclStyle, caption=main.tf]
terraform {
  required_providers {
    docker = {
      source  = "kreuzwerker/docker"
      version = "~> 3.0.2"
    }
  }
}

provider "docker" {}

# --- 1. Base de Donnees PostgreSQL ---
resource "docker_image" "postgres_image" {
  name         = "postgres:latest"
  keep_locally = true
}

resource "docker_container" "db_container" {
  name  = "tp-db-postgres"
  image = docker_image.postgres_image.image_id

  ports {
    internal = 5432
    external = 5432
  }

  env = [
    "POSTGRES_USER=${var.db_user}",
    "POSTGRES_PASSWORD=${var.db_password}",
    "POSTGRES_DB=${var.db_name}",
  ]
}

# --- 2. Application Web Nginx ---
# L'image est pre-construite localement ou via le runner CI/CD
resource "docker_image" "app_image" {
  name = "tp-web-app:latest"
  build {
    context    = "."
    dockerfile = "Dockerfile_app"
  }
  keep_locally = true
}

resource "docker_container" "app_container" {
  name  = "tp-app-web"
  image = docker_image.app_image.image_id

  # La DB doit etre prete avant l'application
  depends_on = [docker_container.db_container]

  ports {
    internal = 80
    external = var.app_port_external
  }
}
\end{lstlisting}


\subsection{outputs.tf}

\begin{lstlisting}[style=hclStyle, caption=outputs.tf]
output "db_container_name" {
  description = "Nom du conteneur de la base de donnees."
  value       = docker_container.db_container.name
}

output "app_access_url" {
  description = "URL d'acces a l'application web."
  value = "http://localhost:${
    docker_container.app_container.ports[0].external
  }"
}
\end{lstlisting}

\section{\textcolor{terraformPurple}{Cycle de Vie du Déploiement (DLC)}}

\subsection{Étape 1 --- Initialisation}

\begin{lstlisting}[style=termStyle, caption=terraform init]
$ terraform init

Initializing the backend...
Initializing provider plugins...
- Finding kreuzwerker/docker versions matching "~> 3.0.2"...
- Installing kreuzwerker/docker v3.0.2...
- Installed kreuzwerker/docker v3.0.2

Terraform has been successfully initialized!
\end{lstlisting}

\begin{tfbox}
\textbf{Rôle :} Télécharge les providers déclarés dans \texttt{required\_providers}
et initialise le répertoire \texttt{.terraform/}. Crée aussi le fichier
\texttt{.terraform.lock.hcl} pour verrouiller les versions et garantir
la reproductibilité.
\end{tfbox}

\subsection{Étape 2 --- Planification}

\begin{lstlisting}[style=termStyle, caption=terraform plan]
$ terraform plan

Terraform will perform the following actions:

  # docker_container.app_container will be created
  + resource "docker_container" "app_container" {
      + name  = "tp-app-web"
      + ports { external = 8080, internal = 80 }
    }

  # docker_container.db_container will be created
  + resource "docker_container" "db_container" {
      + name = "tp-db-postgres"
      + env  = ["POSTGRES_DB=devops_db", ...]
      + ports { external = 5432, internal = 5432 }
    }

  # docker_image.app_image will be created
  + resource "docker_image" "app_image" {
      + name = "tp-web-app:latest"
    }

  # docker_image.postgres_image will be created
  + resource "docker_image" "postgres_image" {
      + name = "postgres:latest"
    }

Plan: 4 to add, 0 to change, 0 to destroy.

Changes to Outputs:
  + app_access_url    = "http://localhost:8080"
  + db_container_name = "tp-db-postgres"
\end{lstlisting}

\begin{tfbox}
\textbf{Rôle :} Simule le déploiement sans modifier quoi que ce soit.
C'est la \textbf{porte d'entrée essentielle} du DLC --- elle permet de valider
les changements avant toute exécution en production.
\end{tfbox}

\subsection{Étape 3 --- Application}

\begin{lstlisting}[style=termStyle, caption=terraform apply]
$ terraform apply -auto-approve

docker_image.postgres_image: Creating...
docker_image.postgres_image: Creation complete after 5s
docker_container.db_container: Creating...
docker_container.db_container: Creation complete after 2s
docker_image.app_image: Creating...
docker_image.app_image: Creation complete after 18s
docker_container.app_container: Creating...
docker_container.app_container: Creation complete after 1s

Apply complete! Resources: 2 added, 0 changed, 0 destroyed.

Outputs:
  app_access_url    = "http://localhost:8080"
  db_container_name = "tp-db-postgres"
\end{lstlisting}

\subsection{Étape 4 --- Validation}

\begin{lstlisting}[style=termStyle, caption=Validation de l'application]
# Verification dans le navigateur
http://localhost:8080
> Application Deployed via Terraform IaC!

# Verification des conteneurs actifs
$ docker ps
CONTAINER ID   IMAGE            STATUS    PORTS
a1b2c3d4e5f6   tp-web-app       Running   0.0.0.0:8080->80/tcp
f6e5d4c3b2a1   postgres:latest  Running   0.0.0.0:5432->5432/tcp
\end{lstlisting}
\begin{figure}[H]
    \centering
    \includegraphics[width=0.5\linewidth]{image.png}
    \label{fig:placeholder}
\end{figure}
\begin{successbox}
\textbf{Résultat :} Les deux conteneurs sont en état \textbf{Running}.
L'application web répond sur le port 8080 et affiche le message de preuve
de déploiement IaC. La base de données PostgreSQL est accessible sur le port 5432.
\end{successbox}

\subsection{Étape 5 --- Destruction}

\begin{lstlisting}[style=termStyle, caption=terraform destroy]
$ terraform destroy -auto-approve

docker_container.app_container: Destroying...
docker_container.app_container: Destruction complete after 1s
docker_container.db_container: Destroying...
docker_container.db_container: Destruction complete after 1s
docker_image.app_image: Destroying...
docker_image.app_image: Destruction complete after 0s
docker_image.postgres_image: Destroying...
docker_image.postgres_image: Destruction complete after 0s

Destroy complete! Resources: 4 destroyed.
\end{lstlisting}

\begin{warnbox}
\textbf{Impact sur terraform.tfstate :} Après \texttt{destroy}, le fichier
\texttt{terraform.tfstate} est mis à jour pour refléter une infrastructure vide.
C'est le mécanisme de \textbf{réconciliation} : Terraform compare l'état réel
avec l'état désiré (vide après destroy) et supprime toutes les ressources.
\end{warnbox}

% ════════════════════════════════════════════
% CI/CD — GITHUB ACTIONS
% ════════════════════════════════════════════
\newpage
\begin{center}
  \rule{\linewidth}{2pt}\\[0.5em]
  {\LARGE\bfseries\textcolor{cicdGreen}{PIPELINE CI/CD}}\\[0.3em]
  {\Large Automatisation du DLC avec GitHub Actions}\\[0.3em]
  \rule{\linewidth}{2pt}
\end{center}

\section{\textcolor{cicdGreen}{Présentation du Pipeline CI/CD}}

L'objectif est d'intégrer les étapes Terraform dans un pipeline afin
d'automatiser le déploiement de l'infrastructure dès qu'une modification
est poussée sur la branche principale.

\begin{cicdbox}
\textbf{Objectifs du Pipeline :}
\begin{itemize}[leftmargin=1.5em]
  \item Pousser le projet sur un dépôt GitHub
  \item Créer le fichier \texttt{.github/workflows/main.yml}
  \item Automatiser Init \textrightarrow{} Plan \textrightarrow{} Apply sur chaque push vers \texttt{main}
  \item Garantir que Apply ne s'exécute que si Plan a réussi
\end{itemize}
\end{cicdbox}

\section{\textcolor{cicdGreen}{Structure du Pipeline}}

\begin{lstlisting}[style=termStyle, caption=Structure du dépôt avec CI/CD]
tp-iac-local/
├── .github/
│   └── workflows/
│       └── main.yml   # Pipeline GitHub Actions
├── main.tf
├── variables.tf
├── outputs.tf
├── Dockerfile_app
├── index.html
├── nginx.conf
└── .gitignore
\end{lstlisting}

\section{\textcolor{cicdGreen}{Fichier de Pipeline --- main.yml}}

\begin{lstlisting}[style=yamlStyle, caption=.github/workflows/main.yml]
name: Terraform CI/CD Pipeline

# Declenchement sur chaque push vers main
on:
  push:
    branches:
      - main

jobs:
  terraform:
    name: Deploy Infrastructure
    runs-on: ubuntu-latest

    steps:
      # Etape 1 : Recupere le code source du depot
      - name: Checkout code
        uses: actions/checkout@v3

      # Etape 2 : Installe Terraform sur le runner
      - name: Setup Terraform
        uses: hashicorp/setup-terraform@v2
        with:
          terraform_version: 1.5.0

      # Etape 3 : Verifie que Docker est disponible
      - name: Verify Docker
        run: docker --version

      # Etape 4 : Initialisation (telecharge le provider Docker)
      - name: Terraform Init
        run: terraform init

      # Etape 5 : Planification (simule les changements)
      - name: Terraform Plan
        run: terraform plan

      # Etape 6 : Application (deploie uniquement si Plan reussi)
      - name: Terraform Apply
        if: success()
        run: terraform apply -auto-approve
\end{lstlisting}

\section{\textcolor{cicdGreen}{Rôle de Chaque Étape dans le DLC}}

\begin{center}
\begin{tabular}{>{\bfseries}p{4cm} p{6cm} p{5cm}}
\toprule
Étape Pipeline & Commande & Rôle dans le DLC \\
\midrule
Checkout        & \texttt{actions/checkout@v3}          & Récupère le code versionné \\
Setup Terraform & \texttt{hashicorp/setup-terraform@v2} & Prépare l'environnement \\
Init            & \texttt{terraform init}                & Initialise l'état \\
Plan            & \texttt{terraform plan}                & Validation \& Revue \\
Apply           & \texttt{terraform apply -auto-approve} & Déploiement Continu \\
\bottomrule
\end{tabular}
\end{center}

\begin{cicdbox}
\textbf{Condition \texttt{if: success()} :} L'étape Apply ne s'exécute que si
toutes les étapes précédentes ont réussi --- notamment le Plan. Cela garantit
qu'aucun déploiement erroné n'est appliqué en cas d'erreur de configuration.
\end{cicdbox}

\section{\textcolor{cicdGreen}{Exécution du Pipeline}}

\subsection{Initialisation du dépôt Git}

\begin{lstlisting}[style=termStyle, caption=Push initial vers GitHub]
$ git init
$ git add .
$ git commit -m "Initial commit - IaC Terraform local"
$ git remote add origin https://github.com/user/tp-iac-local.git
$ git branch -M main
$ git push -u origin main
\end{lstlisting}

\subsection{Résultat du Pipeline sur GitHub Actions}

\begin{lstlisting}[style=termStyle, caption=Logs GitHub Actions - Pipeline reussi]
Run hashicorp/setup-terraform@v2
  Terraform v1.5.0 installed

Run terraform init
  Initializing provider plugins...
  - Installing kreuzwerker/docker v3.0.2...
  Terraform has been successfully initialized!

Run terraform plan
  Plan: 4 to add, 0 to change, 0 to destroy.

Run terraform apply -auto-approve
  docker_image.postgres_image: Creating...
  docker_image.app_image: Creating... (build from Dockerfile_app)
  docker_container.db_container: Creating...
  docker_container.app_container: Creating...
  Apply complete! Resources: 4 added, 0 changed, 0 destroyed.
\end{lstlisting}
\begin{figure}[H]
    \centering
    \includegraphics[width=0.8\linewidth]{image2.png}
    \label{fig:placeholder}
\end{figure}
\begin{successbox}
\textbf{Résultat :} Le pipeline s'est exécuté avec succès. Chaque push vers
\texttt{main} déclenche automatiquement le cycle complet
Init \textrightarrow{} Plan \textrightarrow{} Apply, garantissant un déploiement
continu et reproductible de l'infrastructure.
\end{successbox}

\subsection{Validation de l'Automatisation}

Pour tester le déclenchement automatique, on modifie une variable ,le mot de passe de la DB puis on pousse :

\begin{lstlisting}[style=termStyle, caption=Test de re-deploiement automatique]
# Modification d'une variable
$ # Editer variables.tf : changer db_password

$ git add variables.tf
$ git commit -m "Update: change db_password variable"
$ git push

# GitHub Actions se declenche automatiquement
# Le pipeline re-execute Init -> Plan -> Apply
# La ressource est mise a jour sans intervention manuelle
\end{lstlisting}
\begin{figure}[H]
    \centering
    \includegraphics[width=0.8\linewidth]{image3.png}
    \label{fig:placeholder}
\end{figure}

% ════════════════════════════════════════════
% PARTIE II — PULUMI
% ════════════════════════════════════════════
\newpage
\begin{center}
  \rule{\linewidth}{2pt}\\[0.5em]
  {\LARGE\bfseries\textcolor{pulumiBlue}{PARTIE II}}\\[0.3em]
  {\Large Infrastructure as Code avec Pulumi \& TypeScript}\\[0.3em]
  {\normalsize Option B --- Infrastructure Locale Simulée (LocalStack)}\\[0.3em]
  \rule{\linewidth}{2pt}
\end{center}

\section{\textcolor{pulumiBlue}{Présentation de Pulumi}}

Pulumi est un outil IaC moderne qui permet de définir l'infrastructure dans
un \textbf{vrai langage de programmation} (TypeScript, Python, Go...).
Contrairement à Terraform, il offre toute la puissance du langage :
boucles, conditions, fonctions, types, etc.

\begin{pubox}
\textbf{Objectifs de la Partie II :}
\begin{itemize}[leftmargin=1.5em]
  \item Initialiser un projet Pulumi TypeScript avec \texttt{pulumi new aws-typescript}
  \item Déployer un réseau Docker, PostgreSQL et Nginx
  \item Simuler un bucket S3 AWS via LocalStack
  \item Gérer les secrets avec \texttt{config.getSecret()}
  \item Maîtriser le cycle de vie : \texttt{preview} \textrightarrow{} \texttt{up} \textrightarrow{} \texttt{destroy}
\end{itemize}
\end{pubox}

\section{\textcolor{pulumiBlue}{Initialisation du Projet}}

\subsection{Création du projet}

\begin{lstlisting}[style=termStyle, caption=Initialisation avec pulumi new]
PS> pulumi new aws-typescript
Project name    : tp-pulumi-iac
Project description : TP IaC DevOps
Stack name      : dev
AWS region      : us-east-1

Installing dependencies...
added 327 packages in 2m
Your new project is ready to go!
\end{lstlisting}

\subsection{Structure générée}

\begin{lstlisting}[style=termStyle, caption=Structure du projet Pulumi]
tp3-devops/
├── index.ts          # Code principal (ressources IaC)
├── Dockerfile_app    # Image Nginx personnalisee
├── Pulumi.yaml       # Metadonnees du projet
├── Pulumi.dev.yaml   # Configuration du stack dev
├── package.json      # Dependances Node.js
└── tsconfig.json     # Configuration TypeScript
\end{lstlisting}

\subsection{Installation des dépendances}

\begin{lstlisting}[style=termStyle, caption=Installation des packages Pulumi]
PS> npm install @pulumi/docker
PS> npm install @pulumi/aws
\end{lstlisting}

\subsection{Configuration du stack}

\begin{lstlisting}[style=termStyle, caption=Configuration des variables]
pulumi config set aws:accessKey test
pulumi config set aws:secretKey test
pulumi config set aws:skipCredentialsValidation true
pulumi config set aws:skipRequestingAccountId true
pulumi config set aws:region us-east-1
pulumi config set aws:endpoints '[{"s3":"http://localhost:4566"}]'
pulumi config set tp-pulumi-iac:dbName devops_db
pulumi config set tp-pulumi-iac:dbUser devops_user
pulumi config set --secret tp-pulumi-iac:dbPassword strongpassword123
pulumi config set tp-pulumi-iac:appPortExternal 8080
\end{lstlisting}

Le fichier \texttt{Pulumi.dev.yaml} résultant :

\begin{lstlisting}[style=hclStyle, caption=Pulumi.dev.yaml]
config:
  aws:region: us-east-1
  aws:accessKey: test
  aws:secretKey: test
  aws:skipCredentialsValidation: "true"
  aws:skipRequestingAccountId: "true"
  aws:endpoints: '[{"s3":"http://localhost:4566"}]'
  tp-pulumi-iac:dbName: devops_db
  tp-pulumi-iac:dbUser: devops_user
  tp-pulumi-iac:dbPassword:
    secure: AAABAFuPL/1mXo9jXYQnEYJKexULq740BtK/+TKunCY5...
  tp-pulumi-iac:appPortExternal: "8080"
\end{lstlisting}

\begin{pubox}
\textbf{Sécurité des secrets :} Le mot de passe est automatiquement chiffré
par Pulumi via \texttt{config.getSecret()}. Il est stocké chiffré dans le YAML
et n'apparaît jamais en clair dans les logs.
\end{pubox}

\section{\textcolor{pulumiBlue}{Code Infrastructure --- index.ts}}

\subsection{Variables configurables}

\begin{lstlisting}[style=tsStyle, caption=Variables (équivalent variables.tf)]
// ============================================
// 1. VARIABLES CONFIGURABLES (équiv. variables.tf)
// ============================================
const config = new pulumi.Config();
const dbName = config.get("dbName") || "devops_db";
const dbUser = config.get("dbUser") || "devops_user";
const dbPassword = config.getSecret("dbPassword") || pulumi.secret("strongpassword123");
const appPortExternal = config.getNumber("appPortExternal") || 8080;
const bucketFixedName = config.get("bucketName") || "tp-static-files";

const awsLocalProvider = new aws.Provider("aws-local-provider", {
    region: "us-east-1",
    accessKey: "test",
    secretKey: "test",
    s3UsePathStyle: true,
    skipCredentialsValidation: true,
    skipMetadataApiCheck: true,
    skipRegionValidation: true,
    skipRequestingAccountId: true,
    endpoints: [{ s3: "http://localhost:4566" }],
});
\end{lstlisting}

\subsection{Réseau Docker}

\begin{lstlisting}[style=tsStyle, caption=Création du réseau partagé]
// ============================================
// 2. RESEAU DOCKER
// ============================================
const network = new docker.Network("tp-network", {
    name: "tp-iac-network",
});

\end{lstlisting}

\subsection{Conteneur PostgreSQL}

\begin{lstlisting}[style=tsStyle, caption=Déploiement de la base de données]
// ============================================
// 3. BASE DE DONNEES — Conteneur PostgreSQL
// ============================================
const dbContainer = new docker.Container("db-container", {
    name: "tp-db-postgres",
    image: "postgres:latest",
    networksAdvanced: [{
        name: network.name,
    }],
    ports: [{
        internal: 5432,
        external: 5432,
    }],
    envs: [
        pulumi.interpolate`POSTGRES_USER=${dbUser}`,
        pulumi.interpolate`POSTGRES_PASSWORD=${dbPassword}`,
        pulumi.interpolate`POSTGRES_DB=${dbName}`,
    ],
    restart: "always",
});

\end{lstlisting}

\subsection{Application Web Nginx}

\begin{lstlisting}[style=tsStyle, caption=Build et déploiement de l'application]

// ============================================
// 4. APPLICATION WEB — Conteneur Nginx
// ============================================
const appImage = new docker.Image("app-image", {
    imageName: "tp-web-app:latest",
    build: {
        context: ".",
        dockerfile: "Dockerfile_app",
    },
    skipPush: true,
});

const appContainer = new docker.Container("app-container", {
    name: "tp-app-web",
    image: appImage.imageName,
    networksAdvanced: [{
        name: network.name,
    }],
    ports: [{
        internal: 80,
        external: appPortExternal,
    }],
    restart: "always",
}, {
    dependsOn: [dbContainer],
});


const bucket = new aws.s3.Bucket("static-bucket", {
    bucket: bucketFixedName,
}, {
    provider: awsLocalProvider,
});

const file = new aws.s3.BucketObject("index-html", {
    bucket: bucket.id,
    key: "index.html",
    content: "<h1>Hello from LocalStack S3 ��</h1>",
    contentType: "text/html",
}, {
    provider: awsLocalProvider,
});

\end{lstlisting}

\subsection{Bucket S3 (LocalStack)}

\begin{lstlisting}[style=tsStyle, caption=Bucket S3 simulé via LocalStack]
const awsLocalProvider = new aws.Provider("aws-local-provider", {
    region: "us-east-1",
    accessKey: "test",
    secretKey: "test",
    s3UsePathStyle: true,
    skipCredentialsValidation: true,
    skipMetadataApiCheck: true,
    skipRegionValidation: true,
    skipRequestingAccountId: true,
    endpoints: [{ s3: "http://localhost:4566" }],
});
\end{lstlisting}

\subsection{Outputs}

\begin{lstlisting}[style=tsStyle, caption=Outputs du déploiement]
// ============================================
// 6. OUTPUTS (équiv. outputs.tf)
// ============================================
export const dbContainerName = dbContainer.name;
export const appAccessUrl = pulumi.interpolate`http://localhost:${appPortExternal}`;
export const networkName = network.name;
export const bucketName = bucketFixedName;
export const fileUrl = pulumi.interpolate`http://localhost:4566/${bucketFixedName}/index.html`;
\end{lstlisting}

\section{\textcolor{pulumiBlue}{Cycle de Vie du Déploiement (DLC)}}

\subsection{Étape 1 --- Preview}

\begin{lstlisting}[style=termStyle, caption=pulumi preview]
PS> pulumi preview

     Type                               Name                Plan
 +   pulumi:pulumi:Stack                tp-pulumi-iac-dev   create
 +   |- docker:index:Network            tp-network          create
 +   |- docker:index:Image              app-image           create
 +   |- docker:index:Container          db-container        create
 +   |- docker:index:Container          app-container       create
 +   |- aws:s3:Bucket                   static-bucket       create
 +   |- aws:s3:BucketPublicAccessBlock  bucket-public-...   create
 +   +- aws:s3:BucketObject             index-html          create

Resources: + 8 to create

Outputs:
    appAccessUrl   : "http://localhost:8080"
    dbContainerName: "tp-db-postgres"
    networkName    : "tp-iac-network"
\end{lstlisting}

\subsection{Étape 2 --- Déploiement}

\begin{lstlisting}[style=termStyle, caption=pulumi up]
PS> pulumi up

Updating (dev)

 +  docker:index:Network    tp-network      created  (2s)
 +  docker:index:Image      app-image       created  (31s)
 +  docker:index:Container  db-container    created  (55s)
 +  docker:index:Container  app-container   created  (1s)
 +  aws:s3:Bucket           static-bucket   created
 +  aws:s3:BucketObject     index-html      created
Outputs:
    appAccessUrl   : "http://localhost:8080"
    bucketName     : "tp-static-files"
    dbContainerName: "tp-db-postgres"
    fileUrl        : "http://localhost:4566/tp-static-files/index.html"
    networkName    : "tp-iac-network"
Resources: 8 created
Duration: ~2 minutes
\end{lstlisting}

\subsection{Étape 3 --- Validation}

\begin{lstlisting}[style=termStyle, caption=Validation - conteneurs actifs]
PS> docker ps
CONTAINER ID   IMAGE             STATUS    PORTS
a1b2c3d4e5f6   localstack        Running   0.0.0.0:4566->4566/tcp
3df25a77d7db   postgres:latest   Running   0.0.0.0:5432->5432/tcp
a16ff95714e9   tp-web-app        Running   0.0.0.0:8080->80/tcp

PS> curl http://localhost:8080
<h1>Application Deployed via Pulumi IaC!</h1>

PS> curl http://localhost:4566/tp-static-files/index.html
<h1>Static file hosted via Pulumi IaC!</h1>
\end{lstlisting}
\begin{figure}[H]
    \centering
    \includegraphics[width=0.8\linewidth]{image4.png}
    
    \label{fig:placeholder}
\end{figure}

\begin{successbox}
\textbf{Résultat} : Les 3 conteneurs (\texttt{localstack}, \texttt{tp-db-postgres},
\texttt{tp-app-web}) sont en état \textbf{Running} et toutes les ressources
répondent correctement.
\end{successbox}

\subsection{Étape 4 --- Destruction}

\begin{lstlisting}[style=termStyle, caption=pulumi destroy]
PS> pulumi destroy

Destroying (dev)

 -  aws:s3:BucketObject     index-html      deleted
 -  aws:s3:Bucket           static-bucket   deleted
 -  docker:index:Container  app-container   deleted
 -  docker:index:Container  db-container    deleted
 -  docker:index:Image      app-image       deleted
 -  docker:index:Network    tp-network      deleted
Resources: 8 deleted
\end{lstlisting}
% ════════════════════════════════════════════

% ════════════════════════════════════════════
\section{Questions d'Approfondissement}

\subsection{Impact de terraform destroy sur terraform.tfstate}

Le fichier \texttt{terraform.tfstate} est le \textbf{registre de vérité} de
Terraform : il mappe chaque ressource du code à son identifiant réel dans
l'infrastructure. Lors d'un \texttt{terraform destroy}, Terraform lit ce fichier
pour identifier toutes les ressources à supprimer, les supprime une par une,
puis met à jour \texttt{tfstate} pour refléter une infrastructure vide.
C'est le mécanisme de \textbf{réconciliation} : Terraform compare l'état réel
(ressources existantes) à l'état désiré (vide) et applique les différences.

\subsection{Immuabilité de l'Infrastructure}

L'immuabilité signifie que les ressources \textbf{ne sont jamais modifiées
en place}. Si une configuration change, l'outil IaC détruit l'ancienne ressource
et en crée une nouvelle identique avec la nouvelle configuration. Cela garantit
la \textbf{reproductibilité} et évite la \textit{configuration drift}
(dérive silencieuse de la configuration due à des modifications manuelles).

\subsection{Pourquoi le Plan/Preview est la porte d'entrée du DLC}

L'étape \texttt{plan} (Terraform) ou \texttt{preview} (Pulumi) est
essentielle car elle permet de :
\begin{itemize}
  \item \textbf{Détecter les erreurs} de configuration avant tout déploiement
  \item \textbf{Visualiser exactement} quelles ressources seront créées, modifiées ou détruites
  \item \textbf{Prévenir les suppressions accidentelles} en production
  \item \textbf{Servir de revue de code} pour l'infrastructure dans un pipeline CI/CD
\end{itemize}

\subsection{Alternative : Réseau Docker dans Terraform}

Pour connecter les deux conteneurs dans un réseau Docker partagé avec
Terraform, il faudrait utiliser la ressource \texttt{docker\_network} :

\begin{lstlisting}[style=hclStyle, caption=Réseau Docker dans Terraform]
resource "docker_network" "tp_network" {
  name = "tp-iac-network"
}

# Puis dans chaque conteneur :
resource "docker_container" "app_container" {
  networks_advanced {
    name = docker_network.tp_network.name
  }
}
\end{lstlisting}

% ════════════════════════════════════════════
\section{Comparaison Terraform vs Pulumi}

\begin{center}
\begin{tabular}{>{\bfseries}p{4cm}
    >{\centering\arraybackslash}p{4cm}
    >{\centering\arraybackslash}p{5cm}}
\toprule
Critère &
\textcolor{terraformPurple}{\textbf{Terraform}} &
\textcolor{pulumiBlue}{\textbf{Pulumi}} \\
\midrule
Langage        & HCL (déclaratif)           & TypeScript, Python, Go \\
Courbe apprent.& Faible (HCL simple)        & Moyenne (langage requis) \\
Flexibilité    & Limitée                    & Totale (boucles, conditions) \\
Fichiers       & \texttt{.tf} multiples     & Un seul \texttt{index.ts} \\
État           & \texttt{terraform.tfstate} & Pulumi Cloud (backend) \\
Secrets        & Variables d'env.           & \texttt{config.getSecret()} \\
Preview        & \texttt{terraform plan}    & \texttt{pulumi preview} \\
Déploiement    & \texttt{terraform apply}   & \texttt{pulumi up} \\
CI/CD          & GitHub Actions natif       & GitHub Actions natif \\
Écosystème     & Très mature                & En croissance rapide \\
\bottomrule
\end{tabular}
\end{center}

% ════════════════════════════════════════════
\section{Conclusion}

Ce TP a permis de maîtriser deux approches complémentaires de l'IaC ainsi
que leur automatisation via CI/CD :

\begin{successbox}
\textbf{Partie I --- Terraform} \checkmark
\begin{itemize}[leftmargin=1.5em,topsep=0pt]
  \item Structure de fichiers HCL (\texttt{main.tf}, \texttt{variables.tf}, \texttt{outputs.tf})
  \item Déploiement PostgreSQL + Nginx via le provider Docker
  \item Cycle DLC complet : init \textrightarrow{} plan \textrightarrow{} apply \textrightarrow{} validate \textrightarrow{} destroy
  \item Résolution des problèmes Windows (encodage, config Nginx)
\end{itemize}
\end{successbox}

\vspace{0.5em}

\begin{successbox}
\textbf{Pipeline CI/CD --- GitHub Actions} \checkmark
\begin{itemize}[leftmargin=1.5em,topsep=0pt]
  \item Dépôt GitHub créé et versionné
  \item Pipeline automatique sur chaque push vers \texttt{main}
  \item Étapes Init \textrightarrow{} Plan \textrightarrow{} Apply avec condition \texttt{if: success()}
  \item Validation du re-déploiement automatique après modification
\end{itemize}
\end{successbox}

\vspace{0.5em}

\begin{successbox}
\textbf{Partie II --- Pulumi (TypeScript)} \checkmark
\begin{itemize}[leftmargin=1.5em,topsep=0pt]
  \item Projet initialisé avec \texttt{pulumi new aws-typescript}
  \item Réseau Docker + PostgreSQL + Nginx + Bucket S3 (LocalStack)
  \item Gestion des secrets avec \texttt{config.getSecret()}
  \item Cycle DLC complet : preview \textrightarrow{} up \textrightarrow{} validate \textrightarrow{} destroy
\end{itemize}
\end{successbox}

\vspace{0.5em}

Les trois composantes convergent vers le même objectif --- une infrastructure
\textbf{reproductible, versionnée et automatisable} --- avec Pulumi offrant
plus de puissance grâce aux langages de programmation, et GitHub Actions
garantissant un déploiement continu sans intervention manuelle.

\end{document}